\par\Needspace{6\baselineskip}\medskip
\begin{center}\hypertarget{native-170-article123}{}\pdfbookmark[2]{123.}{article-123}123.\end{center}
In the foregoing we supposed that approximate values of the unknowns $x,y$ were already available from some source. Whenever the approximate dimensions of the entire orbit are known, perhaps derived by earlier calculations from other observations and now to be corrected by new ones, that condition can be met without difficulty, whatever meaning we assign to the unknowns. In the initial determination of an orbit that is still entirely unknown, however, the choice of unknowns matters a great deal. This is by far the most difficult problem. The unknowns must be chosen judiciously so that their approximate values can be obtained from the nature of the problem itself. This succeeds best whenever the three observations used to investigate the orbit cover a heliocentric motion of the celestial body that is not too large. We should therefore always use such observations for the initial determination, which can then be corrected as desired with observations farther apart. It is easy to see that the closer together the observations we use, the more their unavoidable errors disturb the calculation. It follows that the observations for the initial determination must be chosen with care. We must ensure, \emph{first}, that they are not too close together and, \emph{second}, that they are not too far apart. If the observations were too close together, the elements that satisfy them could indeed be calculated with the greatest ease, but little reliance could be placed on the elements themselves. Errors might even distort them so severely that they could not even serve as an approximation. If the observations were too far apart, we would lack the means for an approximate determination of the unknowns. We could then derive \SourceLeaf{171}only a very crude approximation, which would require many more hypotheses, or an entirely unsuitable one, which would make the most tedious trials unavoidable. Judging these limits of the method well is learned better through frequent practice than through rules. The examples given below will show that elements of considerable precision can already be derived from observations of Juno that are only 22 days apart and cover a heliocentric motion of $7^0 35'$. Conversely, our method can still be applied with excellent results to observations of Ceres that are 260 days apart and cover a heliocentric motion of $62^0 55'$. With four hypotheses, or rather four successive approximations, it produces elements that agree very well with the observations.
